% =============================================================================
%  Appendix A — Undermind Group Search Prompts
% =============================================================================

\section{Undermind Group Search Prompts}
\label{app:undermind}

The following six prompts were submitted verbatim to Undermind (undermind.ai) to
retrieve candidate papers for Groups G1--G6 of the pre-validated corpus described in
Section~\ref{sec:methodology:corpora}. Undermind was used exclusively for candidate
\textit{retrieval}; all inclusion and exclusion decisions were made by the author
through manual assessment of each returned record. The prompts are reproduced here
in their entirety for transparency and reproducibility.

% -----------------------------------------------------------------------
\subsection*{Group 1 — Advanced Adapter Architectures \& PEFT Methods}
\label{app:undermind:g1}

\begin{quote}\small
I am researching parameter-efficient fine-tuning (PEFT) methods for large language
models, building on the foundational work of Houlsby et al.\ (2019) on bottleneck
adapters and Hu et al.\ (2022) on LoRA (Low-Rank Adaptation). I am looking for
papers that extend, improve, or propose alternatives to these methods, including but
not limited to: QLoRA (quantized LoRA), DoRA, IA\textsuperscript{3}, prefix tuning,
prompt tuning, BitFit, sparse adapters, hypernetwork-based adapters (HyperFormer),
adapter pruning, rank-adaptive LoRA, VeRA, and methods that combine adapters with
quantization or knowledge distillation. I am specifically interested in works that
analyze the expressiveness, rank selection, or composability of low-rank and bottleneck
adapter modules in transformer architectures (encoder-only, decoder-only, and
encoder-decoder). Please retrieve the most cited and most recent SOTA papers in this
space, including survey papers beyond Han et al.\ (2024, TMLR).
\end{quote}

% -----------------------------------------------------------------------
\subsection*{Group 2 — Multi-Task Learning \& Adapter Composition}
\label{app:undermind:g2}

\begin{quote}\small
I am researching multi-task learning with adapter modules and task composition
techniques for transformer-based language models, building on AdapterFusion (Pfeiffer
et al., 2021, EACL) and AdapterHub (Pfeiffer et al., 2020, EMNLP). I want to find
papers on: non-destructive composition of multiple task-specific adapters, adapter
stacking and fusion strategies, hypernetwork-based multi-task adapter generation
(e.g., HyperFormer, HyperFormer++), multi-task adapter routing, knowledge transfer
between adapters, continual learning with adapters, task arithmetic and model merging
(e.g., TIES-merging, DARE), and zero-shot or few-shot generalization via adapter
composition. I am also interested in works comparing centralized multi-task
fine-tuning with modular adapter-based multi-task approaches. Please retrieve the
most relevant and highly cited SOTA papers, including recent 2023--2025 works.
\end{quote}

% -----------------------------------------------------------------------
\subsection*{Group 3 — P2P Networks \& Decentralised Systems for ML}
\label{app:undermind:g3}

\begin{quote}\small
I am researching peer-to-peer (P2P) and decentralized distributed systems applied to
machine learning and deep learning, building on: Petals (Borzunov et al., 2022/2023,
ACL/NeurIPS) for collaborative LLM inference, \v{S}ajina et al.\ (2024, Future
Generation Computer Systems) on multi-task P2P learning with encoder-only
transformers, and \v{S}ajina (2021, University of Rijeka doctoral thesis) on
peer-to-peer deep learning. I am looking for papers on: decentralized model training
and inference over P2P overlays, gossip-based and DHT-based protocols for model or
parameter sharing, capability-aware peer discovery in ML networks, decentralized
federated learning without a central server, split inference across heterogeneous edge
nodes, bandwidth-aware and latency-aware distributed inference, and decentralized
model marketplaces. Please also include foundational P2P systems papers (Chord,
Kademlia, BitTorrent) that are relevant to designing distributed ML systems. Retrieve
the most relevant SOTA papers from 2018--2025.
\end{quote}

% -----------------------------------------------------------------------
\subsection*{Group 4 — Efficient Inference \& Multi-Adapter Serving}
\label{app:undermind:g4}

\begin{quote}\small
I am researching efficient inference systems for large language models with a focus
on serving multiple adapters (LoRA or bottleneck) simultaneously over a single frozen
base model, building on S-LoRA (Sheng et al., 2024, MLSys) and Petals (Borzunov
et al., 2023). I want to find papers on: multi-tenant LLM serving with adapter
multiplexing, dynamic adapter loading and swapping at inference time, memory-efficient
batching strategies for adapter inference (e.g., unified paging, continuous batching),
KV cache management under adapter heterogeneity, model serving on constrained hardware
(edge devices, commodity GPUs), speculative decoding with adapters, pipeline
parallelism for adapter-augmented models, and systems like vLLM, Punica, AlpaServe,
FlexGen, and EdgeServe. I am also interested in work on adapter caching, prefetching
strategies, and latency-accuracy tradeoffs in multi-adapter inference. Please retrieve
the most relevant and most cited SOTA papers from 2020--2025, prioritizing MLSys,
OSDI, EuroSys, SOSP, and NSDI venues.
\end{quote}

% -----------------------------------------------------------------------
\subsection*{Group 5 — Routing \& Mixture-of-Experts}
\label{app:undermind:g5}

\begin{quote}\small
I am researching routing mechanisms and Mixture-of-Experts (MoE) architectures
applied to transformer models, with a focus on their connection to adapter-based and
modular multi-task learning. My starting point includes AdapterFusion (Pfeiffer et
al., 2021) and S-LoRA (Sheng et al., 2024). I want to find papers on: sparse MoE
routing in transformers (Switch Transformer, GShard, GLaM, Mixtral), expert choice
routing and load balancing strategies (Zhou et al., 2022), soft routing and
differentiable gating mechanisms, Mixture-of-Adapters (MoA) and MoE applied
specifically to PEFT/LoRA modules (e.g., MoLoRA, MixLoRA, Adapter-X by Li et al.\
2024), task-conditioned adapter routing, learned routing without task identity, and
routing in decentralized or federated settings. I am also interested in theoretical
analyses of routing stability, collapse, and capacity. Please retrieve the most
relevant SOTA papers from 2020--2025, including both NLP and systems-focused works.
\end{quote}

% -----------------------------------------------------------------------
\subsection*{Group 6 — Federated Learning \& Privacy-Preserving Adaptation}
\label{app:undermind:g6}

\begin{quote}\small
I am researching federated learning with a focus on privacy-preserving fine-tuning and
adapter-based personalization for transformer models. I am looking for papers that
connect to the proposed research on decentralized adapter systems. Starting from the general
federated learning literature (FedAvg by McMahan et al., 2017; FedProx by Li et al.,
2020), I want papers on: federated fine-tuning of large language models using LoRA or
adapter modules (e.g., FedPEFT, FedAdapter), personalized federated learning with
heterogeneous local models (pFedMMA, per-FedAvg), non-IID data distributions and
their effect on adapter convergence, differential privacy in federated adapter
training, cross-silo vs.\ cross-device federated learning for NLP, medical and
clinical federated learning benchmarks (FLamby, MIMIC), communication-efficient
federated learning with sparse or compressed updates, and split learning with adapter
modules. Please retrieve the most relevant and most cited SOTA papers from 2017--2025,
prioritizing works that explicitly address heterogeneous data or privacy constraints in
adapter or PEFT contexts.
\end{quote}

% =============================================================================
%  Appendix B — Manual Database Search Queries
% =============================================================================

\section{Manual Database Search Queries}
\label{app:db_queries}

Scopus (Elsevier) and Web of Science (Clarivate) were queried manually as part of the
snowballing and enrichment stages described in Section~\ref{sec:methodology:snowballing}.
Ten advanced search queries were executed on each platform. In Scopus, results were
filtered to journal articles and conference proceedings (\texttt{SRCTYPE("j" OR "p")}).
In WoS, Topic-field (\texttt{TS=}) queries were used with equivalent year filters.
The WoS Starter API was subsequently used in a \textit{post-hoc enrichment role} only
--- to append WoS accession numbers and citation counts to already-retrieved records ---
and was not a source of new candidates. All raw result sets from both platforms are
archived in the project repository (\texttt{exports/scopus/} and \texttt{exports/wos/}),
available at the public GitHub repository shared with this submission.

% -----------------------------------------------------------------------
\subsection*{Scopus Advanced Search Queries}
\label{app:db_queries:scopus}

\paragraph{Q1 — Core Adapter-Based NLP.}
\begin{verbatim}
TITLE-ABS-KEY(("adapter" OR "adapters" OR "LoRA" OR "low-rank")
AND ("parameter-efficient" OR "PEFT" OR "parameter efficient")
AND ("NLP" OR "natural language processing" OR "transformer"))
AND SRCTYPE("j" OR "p")
AND PUBYEAR > 2018
\end{verbatim}

\paragraph{Q2 — Multi-Task Learning with Transformers.}
\begin{verbatim}
TITLE-ABS-KEY((("adapter*" OR "LoRA" OR "parameter-efficient fine-tuning")
AND ("multi-task" OR "multitask" OR "multi-adapter")
AND ("inference" OR "serving" OR "deployment" OR "edge")))
AND SRCTYPE(j OR p)
AND PUBYEAR > 2020
AND NOT TITLE-ABS-KEY(("medical" OR "healthcare" OR "clinical"
    OR "sentiment" OR "recommendation"))
\end{verbatim}

\paragraph{Q3 — Peer-to-Peer \& Distributed NLP.}
\begin{verbatim}
TITLE-ABS-KEY(("peer-to-peer" OR "P2P" OR "decentralized" OR "federated")
AND ("learning" OR "training" OR "inference")
AND ("NLP" OR "natural language" OR "machine learning" OR "deep learning"))
AND SRCTYPE("j" OR "p")
AND PUBYEAR > 2017
\end{verbatim}

\paragraph{Q4 — Adapter Fusion \& Dynamic Routing.}
\begin{verbatim}
TITLE-ABS-KEY(("adapter" OR "adapters")
AND ("fusion" OR "routing" OR "dynamic" OR "mixture-of-experts"
    OR "MoE" OR "gating"))
AND SRCTYPE("j" OR "p")
AND PUBYEAR > 2020
\end{verbatim}

\paragraph{Q5 — Parameter-Efficient Multi-Task Transformers.}
\begin{verbatim}
TITLE-ABS-KEY(("parameter-efficient" OR "PEFT" OR "efficient adaptation")
AND ("multi-task" OR "multitask")
AND ("transformer" OR "pre-trained models"))
AND SRCTYPE("j" OR "p")
AND PUBYEAR > 2019
\end{verbatim}

\paragraph{Q6 — Modular Inference \& Task Switching.}
\begin{verbatim}
TITLE-ABS-KEY(("modular" OR "modularity")
AND ("inference" OR "inference time")
AND ("adapter" OR "task-specific" OR "lightweight modules"))
AND SRCTYPE("j" OR "p")
AND PUBYEAR > 2018
\end{verbatim}

\paragraph{Q7 — Agent-Based Collaborative Learning.}
\begin{verbatim}
TITLE-ABS-KEY(("agent" OR "agents" OR "node" OR "nodes")
AND ("collaborative" OR "collaboration" OR "decentralized")
AND ("deep learning" OR "neural network" OR "model training"))
AND SRCTYPE("j" OR "p")
AND PUBYEAR > 2017
\end{verbatim}

\paragraph{Q8 — NL-to-SQL with Adapters.}
\begin{verbatim}
TITLE-ABS-KEY(("NL-to-SQL" OR "semantic parsing"
    OR "natural language to SQL")
AND ("adapter" OR "fine-tun*" OR "parameter-efficient" OR "efficient"))
AND SRCTYPE("j" OR "p")
AND PUBYEAR > 2019
\end{verbatim}

\paragraph{Q9 — Sequence Labelling \& Token Classification with PEFT.}
\begin{verbatim}
TITLE-ABS-KEY(("token classification" OR "NER"
    OR "named entity recognition" OR "sequence labeling")
AND ("adapter" OR "LoRA" OR "parameter-efficient" OR "PEFT"))
AND SRCTYPE("j" OR "p")
AND PUBYEAR > 2019
\end{verbatim}

\paragraph{Q10 — Frozen Base Model Paradigm.}
\begin{verbatim}
TITLE-ABS-KEY(("frozen" OR "freeze" OR "frozen base")
AND ("adapter" OR "efficient" OR "modular")
AND ("transfer learning" OR "fine-tun*"))
AND SRCTYPE("j" OR "p")
AND PUBYEAR > 2018
\end{verbatim}

% -----------------------------------------------------------------------
\subsection*{Web of Science Advanced Search Queries}
\label{app:db_queries:wos}

\paragraph{Q1 — Core Adapter-Based NLP.}
\begin{verbatim}
TS=(("adapter" OR "adapters" OR "LoRA" OR "low-rank")
AND ("parameter-efficient" OR "PEFT" OR "parameter efficient")
AND ("NLP" OR "natural language processing" OR "transformer"))
AND PY >= 2018
\end{verbatim}

\paragraph{Q2 — Multi-Task Adapter Inference.}
\begin{verbatim}
TS=((("adapter*" OR "LoRA" OR "PEFT")
AND ("multi-task" OR "multitask")
AND ("inference" OR "serving" OR "edge" OR "distributed")))
AND PY >= 2021
AND NOT TS=(("medical" OR "healthcare" OR "sentiment"))
\end{verbatim}

\paragraph{Q3 — Peer-to-Peer Distributed Learning.}
\begin{verbatim}
TS=(("peer-to-peer" OR "P2P" OR "decentralized" OR "federated")
AND ("learning" OR "training" OR "inference")
AND ("NLP" OR "natural language" OR "machine learning"
    OR "neural network"))
AND PY >= 2017
\end{verbatim}

\paragraph{Q4 — Adapter Fusion \& Routing Mechanisms.}
\begin{verbatim}
TS=(("adapter" OR "adapters")
AND ("fusion" OR "dynamic routing" OR "mixture-of-experts"
    OR "MoE" OR "gating"))
AND PY >= 2020
\end{verbatim}

\paragraph{Q5 — Modular Deep Learning.}
\begin{verbatim}
TS=(("modular" OR "modularity")
AND ("neural" OR "transformer" OR "deep learning")
AND ("inference" OR "task-specific" OR "lightweight"))
AND PY >= 2018
\end{verbatim}

\paragraph{Q6 — Agent-Based Collaborative ML.}
\begin{verbatim}
TS=(("agent" OR "agents" OR "distributed node*")
AND ("collaborative" OR "collaboration" OR "decentralized"
    OR "gossip")
AND ("deep learning" OR "machine learning"))
AND PY >= 2017
\end{verbatim}

\paragraph{Q7 — Memory-Efficient Multi-Task NLP.}
\begin{verbatim}
TS=(("multi-task" OR "multitask")
AND ("memory-efficient" OR "memory efficient"
    OR "parameter-efficient" OR "computational efficient")
AND ("NLP" OR "natural language" OR "transformer"))
AND PY >= 2018
\end{verbatim}

\paragraph{Q8 — Transfer Learning with Task Adaptation.}
\begin{verbatim}
TS=(("transfer learning" OR "fine-tun*")
AND ("adapter" OR "parameter-efficient" OR "LoRA")
AND ("NLP" OR "text" OR "language"))
AND PY >= 2019
\end{verbatim}

\paragraph{Q9 — Sequence Task Learning.}
\begin{verbatim}
TS=(("token classification" OR "NER" OR "named entity"
    OR "sequence labeling")
AND ("adapter" OR "fine-tun*" OR "efficient" OR "modular"))
AND PY >= 2019
\end{verbatim}

\paragraph{Q10 — Distributed Inference.}
\begin{verbatim}
TS=(("distributed" OR "decentralized")
AND ("inference" OR "inference time")
AND ("efficient" OR "modular" OR "lightweight"))
AND PY >= 2018
\end{verbatim}

% =============================================================================
%  Appendix C — Search Execution Log
% =============================================================================

\section{Search Execution Log}
\label{app:search_log}

This appendix provides a chronological account of every search activity performed
during corpus construction, together with record counts at each stage. The purpose
is to allow full reproduction or audit of the PRISMA flow by the reader or supervisor.
All output files referenced below are committed to the project repository under
\texttt{slr\_engine/snowballing/snowball\_output/}.

% -----------------------------------------------------------------------
\subsection*{C.1 — Search Timeline}
\label{app:search_log:timeline}


\begin{table}[H]
\centering
\small
\caption{Chronological search log with input and output record counts.}
\label{tab:search_log}
\begin{tabularx}{\textwidth}{@{}
  l
  X
  >{\raggedright}p{2.2cm}
  >{\raggedright\arraybackslash}p{3.8cm}
@{}}
\toprule
\textbf{Date} & \textbf{Activity} & \textbf{In} & \textbf{Out} \\
\midrule
2026-04-21 & Semantic Scholar API snowball on G0 seeds (backward + forward) &
  9 seeds & 1\,150 raw \\
2026-04-21 & Within-snowball deduplication (DOI + SS paper ID) &
  1\,150 & 972 \\
2026-04-21 & Title screening (Layer~1: year; Layer~2: keyword; Layer~3: LLM triage) &
  972 & 162 INCL, 791 EXCL, 19 UNCERT \\
2026-04-21 & Merge snowball INCLUDE with 352-paper pre-validated corpus (G0--G6);
  cross-deduplication & 514 & 502 \\
$\approx$2026-04-22 & Semantic Scholar / OpenAlex enrichment + topical-relevance filter &
  502 & 464 kept; 32 depri.; 6 excl. \\
$\approx$2026-04-28 & Abstract-level review (AI-assisted, human decision) &
  464 & 214 KEEP; 338 SKIP \\
2026-05-02  & Fulltext review queue compiled (KEEP + borderline DEFER) &
  214 KEEP & 387 queued \\
2026-05-08  & Forward citation snowball: Scopus (citing G0) &
  106 raw & 74 title-screened in \\
2026-05-08  & Forward citation snowball: Web of Science (citing G0) &
  20 raw  & 14 title-screened in \\
2026-05-08  & Forward snowball records added to abstract review (G\_SNOW\_F group) &
  88 added & 12 KEEP; 76 SKIP \\
2026-05-11  & DOI-based venue validation: Scopus (\texttt{M3\_scopus\_venue}) &
  172 DOIs & 172 matched \\
2026-05-11  & DOI-based venue validation: Web of Science (\texttt{M4\_wos\_venue}) &
  125 DOIs & 125 matched \\
2026-05-12  & Full-text reading, quality assessment, and final curation &
  387 queued & 123 final \\
\bottomrule
\end{tabularx}
\end{table}

% -----------------------------------------------------------------------
\subsection*{C.2 — Automated Citation Snowballing}
\label{app:search_log:snowball}

Citation snowballing was executed programmatically on \textbf{2026-04-21} via the
Semantic Scholar Academic Graph API, supplemented by Scopus and ACL Anthology engines,
following the procedure of \citet{Wohlin2014}. The eight G0 seed papers were submitted
for both backward (reference list) and forward (citing papers) retrieval. Raw API
responses yielded \textbf{1,150 candidate records}; after deduplication by DOI and
Semantic Scholar paper identifier, \textbf{972 unique records} entered title screening.
Screening decisions are logged in \texttt{log\_screening\_2026-04-21.json}; per-seed
retrieval counts are in \texttt{log\_retrieval\_2026-04-21.json}.

% -----------------------------------------------------------------------
\subsection*{C.3 — Forward Citation Snowballing (Scopus and WoS)}
\label{app:search_log:forward}

On \textbf{2026-05-08}, a supplementary forward snowball was conducted by manually
querying the Scopus and Web of Science interfaces for papers citing any of the G0 seed
papers. This step retrieved \textbf{106 records from Scopus}
(\texttt{MANUAL\_FORWARD\_SCOPUS\_0805.csv}) and \textbf{20 records from WoS}
(\texttt{M1\_forward\_wos\_2026-05-08.txt}). After manual title screening, \textbf{88
records} (74 Scopus + 14 WoS) passed the relevance threshold and were appended to the
abstract review pool under the source group \texttt{G\_SNOW\_F}. These 88 records
bypassed the automated enrichment filter and entered the review directly. Of the 88,
\textbf{12 were assigned KEEP} and 76 were SKIP at abstract review.

% -----------------------------------------------------------------------
\subsection*{C.4 — Manual Database Cross-Validation}
\label{app:search_log:crossval}

To verify completeness, the ten WoS and ten Scopus advanced queries documented in
Appendix~B were also executed manually against the live database interfaces.
\textbf{Scopus} returned \textbf{172 records}
(\texttt{11\_MANUAL\_SCOPUS\_export\_1105.csv}); \textbf{Web of Science} returned
\textbf{125 records} (\texttt{11\_MANUAL\_WOS\_export\_1105.txt}).

Cross-referencing the 125 WoS records against the existing merged corpus revealed that
all candidates were already present via the Semantic Scholar and Scopus retrieval
pipeline; \textbf{zero new candidates} were identified from the WoS manual search.
This outcome is consistent with the known indexing overlap between WoS and Scopus for
the PEFT, NLP systems, and distributed ML literature: Scopus coverage is broader for
conference proceedings from ACL, NeurIPS, ICML, and MLSys venues, which dominate this
review's citation landscape. The WoS manual search therefore served as a \textit{coverage
cross-check} rather than a source of new records, confirming that the Scopus-dominated
retrieval pipeline achieved comprehensive coverage of the target literature.

% -----------------------------------------------------------------------
\subsection*{C.5 — DOI-Based Venue Verification}
\label{app:search_log:doi}

Following full-text reading, the DOIs of candidate papers were submitted in batch to
Scopus and WoS to retrieve authoritative venue, publication year, and citation-count
metadata. The Scopus batch returned \textbf{172 records}
(\texttt{M3\_scopus\_venue\_validation\_2026-05-11.csv}); the WoS batch returned
\textbf{125 records} (\texttt{M4\_wos\_venue\_validation\_2026-05-11.txt}). These
files served exclusively as metadata enrichment sources: they appended WoS accession
numbers and confirmed Scopus EIDs to records already in the reading pool, and were not
used to introduce new candidates.

% -----------------------------------------------------------------------
\subsection*{C.6 — Manual PDF and Venue Check}
\label{app:search_log:pdf}

A subset of candidate papers had been initially downloaded from the arXiv preprint
server. As a final quality step, each arXiv download was manually checked against
Google Scholar, Semantic Scholar, and the ACL Anthology to determine whether a
peer-reviewed published version existed. Where a published venue was identified and
the PDF was publicly accessible, the preprint was replaced with the official published
version and the venue metadata was updated accordingly. After this step,
\textbf{18 papers} in the final 123-paper reading list retained an arXiv source
(\texttt{source = ARXIV} in \texttt{13\_final\_reading\_list\_2026-05-12.csv}),
either because no published version was found or because the accepted manuscript was
not separately accessible. A further \textbf{4 records} have no source annotation,
reflecting papers added from curated references whose primary retrieval provenance was
not logged in the pipeline output files.
